\chapter{Mar.~2 --- Verma Modules}

\section{Integral Form of Weyl Character Formula}

\begin{theorem}
  If $f$ is a conjugation invariant
  continuous function on a compact connected
  Lie group $K$ with maximal torus
  $T \subseteq K$ and Haar probability
  measures $dk$ and $dt$, then
  \[
    \int_K f(k)\, dk
    = \frac{1}{|W|} \int_T f(t) |\delta(t)|^2\, dt,
  \] 
  where
  $\delta(t) = \rho(t)^{-1} \prod_{\alpha \in R_+} (\alpha(t) - 1)$.\footnote{Recall that if $\alpha \in R_+$ and $e^H = t$ for $H \in \mathfrak{t} = \Lie T$, then we denote $e^{\alpha(H)} = t^\alpha = \alpha(t)$.}
\end{theorem}

\begin{proof}
  First (exercise) use
  the Peter-Weyl theorem to show that the
  characters form a dense set in the
  space of conjugation invariant functions.
  So it suffices to show the formula for $f$
  a character. Then
  \[
    \int_K \chi_\lambda(k)\, dk
    = \delta_{\lambda, 0}
  \]
  is the left-hand side
  by the orthogonality of characters.
  Now write
  \[
    \chi_\lambda(t)
    = \frac{\sum_{w \in W} \varepsilon(w) t^{w(\lambda + \rho)}}{\Delta}, \quad
    \Delta = \sum_{w \in W} \varepsilon(w) t^{w(\rho)} = t^{-\rho} \prod_{\alpha \in R_+} (t^\alpha - 1)
    = \rho(t)^{-1} \prod_{\alpha \in R_+} (\alpha(t) - 1),
  \]
  so $\Delta = \delta(t)$.
  Here $\varepsilon(w) = (-1)^{\ell(w)}$
  is the sign of $w$.
  Then we get
  \[
    \frac{1}{|W|} \int_T \chi_\lambda(t)
    \overline{\chi_\mu(t)} |\delta(t)|^2\, dt
    = \frac{1}{|W|}  \int_T
    \sum_{w \in W} \varepsilon(w) t^{w(\lambda + \rho)}
    \sum_{w' \in W} \varepsilon(w') t^{-w'(\mu + \rho)}\, dt
  \]
  This is equal to $\delta_{\lambda, \mu}$
  by the orthogonality of characters, so
  we get the formula.
\end{proof}

\begin{remark}
  If $T \cong \R^r / \Lambda$ where
  $\Lambda \cong \Z^r$ and
  $\mathcal{F}$ is a fundamental domain for
  $T$, then $H = \sum_{i = 1}^r x_i e_i$
  and
  \[
    dt = \frac{dx_1 \ldots dx_r}{\mathrm{vol}(\mathcal{F})}.
  \]
  For $e^{i\theta_j}$ on $U(n)$, this is
  given by $(d\theta_1 \ldots d\theta_n) / (2\pi)^n$.
\end{remark}

\begin{exercise}
  Let $f$ be a conjugation invariant
  function on $U(n)$. Show that
  \[
    \int_{U(n)} f(k)\, dk
    = \frac{1}{(2\pi)^n n!}
    \int_{|z_1| = \cdots = |z_n| = 1}
    f(\diag(z_1, \dots, z_n))
    \prod_{m < j} |z_m - z_j|^2\, d\theta_1 \ldots d\theta_n,
    \quad z_j = e^{i \theta_j}.
  \]
\end{exercise}

\section{More on Verma Modules}

\begin{remark}
  Let $\g$ be a complex semisimple Lie
  algebra and $\mathfrak{h} \subseteq \g$ a Cartan
  subalgebra. Let $\lambda \in \mathfrak{h}^*$
  and $M_\lambda$ the Verma module of
  highest weight $\lambda$, i.e. it is
  generated by one vector $v_\lambda$ with
  \[
    e_i v_\lambda = 0 \quad \text{and} \quad
    h(v_\lambda) = \lambda(h) v_\lambda
  \]
  for $1 \le i \le r$. Note that
  $M_\lambda = \Ind_{\mathfrak{h} \oplus \mathfrak{n}^+} \C_\lambda$,
  where $\C_\lambda$ is the one-dimensional
  representation for
  $\mathfrak{h} \oplus \mathfrak{n}^+$
  defined by by the above actions.
  Note that $M_{\lambda}$ is irreducible
  for generic $\lambda$. In these
  cases,
  \[
    L_\lambda = M_\lambda / I_\lambda
  \]
  is irreducible, where $I_\lambda$ is
  the maximal proper submodule of
  $M_\lambda$. Note that exchanging
  $e_i$ and $f_i$ in the above definitions
  gives the lowest weight Verma modules.

  There is a unique symmetric bilinear
  form on $M_\lambda$ satisfying the
  contravariance property:
  \[
    \langle v_\lambda, v_\lambda = 1 \rangle
    \quad \text{and} \quad
    \langle e_i u, w \rangle 
    = \langle u, f_i w \rangle
  \]
  for $u, w \in M_\lambda$. This
  form is known as the \emph{Shapovalov form},
  and its kernel is precisely $I_\lambda$.
  Note that if $\omega$ is the
  \emph{Chevalley involution} given by
  $\omega(e_i) = -f_i$,
  $\omega(f_i) = -e_i$, and
  $\omega(h_i) = -h_i$, then
  \[
    \langle x u, v \rangle
    = - \langle u, \omega(x) v \rangle
  \]
  is the contravariance property from
  above.

  Let $V = \bigoplus V^\mu$ be the
  weight decomposition (note that the
  weight spaces are finite-dimensional).
  Then the $V^\mu$ and $V^\nu$ are
  orthogonal with respect to the
  Shapovalov form for $\mu \ne \nu$. Let
  \[
    V^* = \bigoplus_{\lambda}\, (V^\lambda)^*
  \]
  be the \emph{restricted dual} of $V$, and
  let $\Hom_{\C}(V_1, V_2)$ be the space
  of all linear maps.
\end{remark}

\begin{prop}
  If $M_\lambda$ is irreducible, then
  its restricted dual is the lowest weight
  Verma module with lowest
  weight $-\lambda$.
\end{prop}

\begin{definition}
  Define
  $V^c$ to be the restricted dual of $V$
  with action of $\g$ twisted by $\omega$:
  \[
    \langle xu^c, v \rangle
    = -\langle u^c, \omega(x) v \rangle,
    \quad u^c \in V^c,\, v \in V,\, x \in \g.
  \]
  Naturally, $V^c$
  has the same weight subspaces as
  $V$. We call
  $M_\lambda^c$ the \emph{contragredient
  Verma module}.
\end{definition}

\begin{remark}
  There is a unique
  $M_\lambda \to M_\lambda^c$ with image
  $L_\lambda$, and
  combining it with the pairing above gives
  \[
    M^c_\lambda \otimes M_\lambda \longrightarrow \C,
  \]
  which is exactly the Shapovalov form.
\end{remark}

\begin{theorem}
  The Verma module is irreducible if
  and only if the following condition is
  satisfied: For any $\alpha \in R_+$,
  we have
  $\langle \lambda + \rho, \alpha^\vee \rangle \notin \N$.
\end{theorem}

\begin{definition}
  We will call $\lambda$ \emph{generic}
  if for all $\alpha \in R_+$,
  we have $\langle \lambda, \alpha^\vee \rangle \notin \Z$.
\end{definition}

\begin{theorem}
  We have the following:
  \begin{enumerate}
    \item The module $L_\lambda$
      is irreducible if and only if
      $\lambda \in P^+$.
    \item The representations $L_\lambda$
      for $\lambda \in P^+$ are pairwise
      non-isomorphic, and any finite-dimensional
      irrep
      is isomorphic to one of them.
    \item If $\lambda = \sum n_i \omega_i$
      with $n_i \in \Z$, then
      $L_\lambda$ is generated by
      $v_\lambda$ and the action of
      $\g$ on $v_\lambda$ with the
      constraints
      $f^{n_i + 1}_i v_\lambda = 0$ for
      $i = 1, \dots, r$.
      Moreover, we have
      $L_\lambda^* \cong L_{-w_0(\lambda)}$.
  \end{enumerate}
\end{theorem}

\begin{prop}
  The category $C(\g)$ of finite-dimensional
  representations is semisimple, i.e.
  every module $V$ is isomorphic to a
  direct sum of simple modules
  $L_\lambda$ for $\lambda \in P^+$.
  Moreover, the map
  \[
    \Phi \longmapsto \Phi(v_\lambda)
  \]
  identifies the multiplicity space
  $\Hom_\g(L_\lambda, V)$ for $\lambda \in P^+$
  with the space
  \[
    (V^\lambda)^{\mathfrak{n}^+}
    = \{v \in V^\lambda : e_i v = 0\}
  \]
  of singular vectors of weight $\lambda$.
\end{prop}

\section{Maximal Roots and the Casimir Operator}

\begin{remark}
  Recall that the \emph{maximal root} $\theta$
  is the highest weight in the
  adjoint representation.
\end{remark}

\begin{prop}
  The maximal root is uniquely determined
  by the following conditions:
  For all $\alpha \in R$, we have
  $\theta - \alpha \in Q^+$.
\end{prop}

\begin{remark}
  We normalize the pairing so that
  $(\theta, \theta) = 2$. Let
  \begin{itemize}
    \item $\theta = \sum_i h_i \alpha_i$,
    \item $h = 1 + \sum_i h_i$,
    \item $h_i^\vee = h_i (\alpha_i, \alpha_i) / 2$ (these are the coefficients of the
      expansion of ${\widetilde{\theta}}^\vee$),
    \item $h^\vee = 1 + \sum_i h_i^\vee = (\rho, \theta) + 1$
      (this is the dual Coxeter number).
  \end{itemize}
  For root systems of type A, D, E, we have
  $h = h^\vee$.
\end{remark}

\begin{remark}
  Recall that $C = \sum x_i x^i$ is the
  \emph{Casimir operator}, where
  $\{x^i\}$ is the dual basis for $x_i$.
  Then
  \[
    \Omega = \sum (x_i \otimes x^i)
  \]
  does not depend on the choice of basis,
  and the value of $C$ on any irreducible
  highest weight module $L_\lambda$ is given by
  $(\lambda, \lambda + 2\rho)$ (and
  for the lowest weight $L_{-\mu}$ by
  $(\mu, \mu + 2\rho)$).
  The value of $C$ in the adjoint representation
  is given by $2h^\vee$.
\end{remark}

\begin{exercise}
  Show that in a tensor product of
  two representations, the action of $\g$
  commutes with $\Omega$.
\end{exercise}

\begin{prop}
  We have
  \[
    \Omega|_{V_\lambda \subseteq V_\mu \otimes V_\nu}
    = \frac{(\lambda, \lambda + 2\rho) - (\mu, \mu + 2\rho) - (\nu, \nu + 2\rho)}{2}.
  \]
\end{prop}

\begin{proof}
  The proof relies on the identity
  \[
    \Omega = \frac{1}{2}(\Delta(C) - c \otimes 1 - 1 \otimes c),
  \]
  where $\Delta : U(\g) \to U(\g) \otimes U(\g)$
  is given by $x \mapsto x \otimes 1 + 1 \otimes x$. Note that
  $(\Delta \otimes I) \Delta = (I \otimes \Delta) \Delta$.
\end{proof}

\section{Affine Lie Algebras}

\begin{remark}
  Consider a manifold $M$ and a 1-form
  $J$ satisfying $d(* J) = 0$. The
  $1$-form $J$ takes values in a Lie
  algebra, and we have
  \[
    [J^a(x), J^b(y)]
    = f^{a, b, c} J^c(x) \delta(x - y),
  \]
  where the $f^{a, b, c}$ are the
  structure constants of the Lie algebra.
  These are called \emph{current algebras},
  see the book by J. Mickellson for
  more details (e.g. for maps $S^3 \to \g$).

  We will consider the case where
  $M$ is 2-dimensional. Then the
  conservation property becomes
  \[
    \partial_{\overline{z}} J_z
    + \partial_z J_{\overline{z}} = 0.
  \]
  In some cases, this can be separated,
  so that $\partial_{\overline{z}} J_z = 0$,
  i.e. $J_z$ is holomorphic. Then we can write
  \[
    J_z = \sum_{n \in \Z} J_n z^{-n - 1}.
  \]
  The $J_n$ in this expansion are called
  \emph{modes} (one can think of this as
  some kind of wave expansion). Then
  the commutation relations become
  (with some correction term)
  \[
    [J^a_n, J^b_m]
    = f^{a, b, c} J^c_{n + m}.
  \]
\end{remark}

\begin{remark}
  One can think of the modes above as
  loops. The \emph{loop algebra} is
  \[
    L \g = \g \otimes \C[t, t^{-1}].
  \]
  One can also think of them as maps
  $S^1 \to \g$. The commutation relations
  for $L \g$ are
  \[
    [x \otimes P, y \otimes Q]
    = [x, y] \otimes PQ,
  \]
  so that $x \otimes t^n = x[n]$ has
  commutation relations
  \[
    [x[n], y[m]] = [x, y][m + n].
  \]
  This algebra has a unique nontrivial central
  extension, i.e. a short exact sequence of
  Lie algebras
  \[
    \begin{tikzcd}
      0 \ar[r] & \g \ar[r] & \widehat{\g}
      \ar[r] & \C c \ar[r] & 0
    \end{tikzcd}
  \]
  Here $\g$ is a simple Lie algebra. Then
  on $\widehat{g}$, we have
  \[
    [x \otimes P, y \otimes Q]
    = [x, y] \otimes PQ
    + c \frac{(x, y)}{2\pi i} \oint_{|t| = 1}Q\, dP,
  \]
  which gives the commutation relation
  \[
    [x[n], y[m]] = [x, y][m + n]
    + n \delta_{n, -m} (x, y) c.
  \]
  Note that we can also extend
  $\widehat{\g}$ by the grading
  $d = t \frac{d}{dt}$ and get
  $\widetilde{g} = \widehat{g} \oplus \C d$.
\end{remark}
